\begin{answers}{ch:modern-arch}{現代の LLM アーキテクチャ}
\item どちらも $(1.265,\,-1.265,\,0.632,\,-0.632)$ で、一致します。
2 乗の平均は $(4+4+1+1)/4=2.5$ なので、$\mathrm{RMS}(\vx)=\sqrt{2.5}\approx1.581$ です。
RMSNorm は各成分をこれで割るので、$2/1.581=1.265$、$1/1.581=0.632$ となります。
LayerNorm では、平均が $\mu=(2-2+1-1)/4=0$ なので、引いても $\vx$ は変わりません。
分散は $\sigma^2=2.5$ で、標準偏差は RMS と同じ $1.581$ です。
平均が 0 のベクトルでは、式\eqref{eq:ma-layernorm}と式\eqref{eq:ma-rmsnorm}は同じ計算になります。
\item どちらも $\sin1\approx0.841$ で、一致します。
位置 $m$ のクエリは $\mR(m\omega)\vq=(\cos m\omega,\,\sin m\omega)$、位置 $n$ のキーは $\mR(n\omega)\vk=(-\sin n\omega,\,\cos n\omega)$ です。
$m=4,\,n=2$ では、クエリが $(\cos2,\sin2)=(-0.416,\,0.909)$、キーが $(-\sin1,\cos1)=(-0.841,\,0.540)$ です。
内積は $0.350+0.491=0.841$ です。
$m=10,\,n=8$ では、クエリが $(0.284,\,-0.959)$、キーが $(0.757,\,-0.654)$ で、内積は $0.215+0.627=0.841$ です。
式\eqref{eq:ma-rope-rel}のとおり、スコアは差 $n-m=-2$ だけで決まり、$\vq^\T\mR(-2\omega)\vk=\sin1$ となるからです。
\item $\dff=\frac{8}{3}\times4096\approx10{,}923$ です。Llama 2 7B の 11008 は、これより約 0.8\% 大きいだけです。
GPT-2 型のパラメータ数は $8d^2\approx1.342$ 億、SwiGLU は $3d\,\dff$ なので、等しいと置くと $\dff=\frac83d=10922.7$ です。
11008 は $43\times256$ で、$10922.7$ を 256 の倍数に切り上げた値になっています。
Llama 2 は、GPT-2 型とほぼ同じパラメータ数（約 1.35 億）になるように $\dff$ を選んでいると分かります。
\item キャッシュのメモリはちょうど 4 分の 1 になり、速さはほとんど変わりませんでした（実測）。
80 トークンを入れた後、1 層の \code{KVCache} の \code{k} と \code{v} の形は、どちらも $(1,4,80,48)$ から $(1,1,80,48)$ になりました。
全 4 層で 1 トークンあたりに保存する数は、1,536 個から 384 個に減ります。
新トークン 480 個の生成（キャッシュあり、3 回の中央値）は、同じ計算機で続けて 2 回ずつ測ると、\code{n_kv_head=4} が 1,183 ms と 1,186 ms、1 が 1,126 ms と 1,111 ms でした（本文の \file{ch11/out/kv_cache.txt} は 1,171 ms）。
速くなったのは 5\% ほどで、メモリが 4 分の 1 になったのに比べると、わずかです。
MQA や GQA の主な利点は、速さよりもメモリにあります。
\item Llama 3.1 405B は 1 トークンあたり 504 KiB、文脈いっぱいで 63 GiB です。gpt-oss-120b は 72 KiB、9 GiB で、7 分の 1 です。
式\eqref{eq:ma-kv-mem}に入れると、405B は $2\times126\times8\times128\times2=516{,}096$ バイトで、1024 で割ると 504 KiB です。
131,072 トークンでは $504\times131{,}072$ KiB $=63$ GiB になります。
gpt-oss-120b は $2\times36\times8\times64\times2=73{,}728$ バイト $=72$ KiB で、131,072 トークンでは 9 GiB です。
層の数が $126/36=3.5$ 倍、$\dhead$ が 2 倍違うので、合わせて 7 倍の差になります。
なお、gpt-oss は一部の層でスライディングウィンドウ注意を使うので、実際のキャッシュは 9 GiB より小さくなります。
\item 検証損失は 4.896 から 4.638 に下がり、エキスパートの割合は均等に少し近づきました（実測、1 回の実行）。FFN の計算量は約 2 倍です。
\code{top_k=2} に変えた MoE（補助損失 0.01）を、同じ 800 ステップで学習させました。
選ばれた割合は、層 0 が $0.31,\,0.20,\,0.20,\,0.30$、層 1 が $0.21,\,0.19,\,0.31,\,0.29$ です。
top-1 では最小が 0.13 でしたが、top-2 では 0.19 になりました。
1 トークンが 2 個のエキスパートを通るので、計算量も、各エキスパートが見るトークンも 2 倍に増えます。
そのため、短い学習でもエキスパートの学習が進みやすくなったと考えられます。
もう一つの違いは、ルータの学習です。本章のコードの top-1 では、割り直した重みが常に 1 になるので、ルータには $\calL_{\mathrm{LM}}$ の勾配がほぼ届かず、補助損失だけで学習していました（\secref{sec:ma-moe}）。top-2 では重みが 2 個の確率の比で決まるので、ルータが $\calL_{\mathrm{LM}}$ からも学習するようになります。
ただし、計算量を 2 倍にした結果なので、密なモデル（4.793）との公平な比較ではありません。
\item GPU のメモリが 24 GB 程度なら、密な 80 億パラメータのモデルが向いています。メモリが 60 GB を超えるなら、MoE のほうが速く動きます。
BF16 の重みは、密なモデルが $80\text{ 億}\times2$ バイト $=16$ GB、MoE は総パラメータ分の $300\text{ 億}\times2=60$ GB です。
1 トークンの計算量はアクティブパラメータで決まるので、MoE は 30 億分で、密なモデルの約 $3/8$ です。
MoE は「メモリを多く使う代わりに計算を減らす」設計なので、載りさえすれば速く生成できます。
メモリが足りない GPU 1 枚では密なモデルを選び、MoE を使いたいときは量子化（\chapref{ch:distillation}）で重みを小さくすることを考えます。
\end{answers}
